% !TeX program = xelatex
\documentclass[a4paper,fontsize=11pt,parskip=half,headings=small]{scrartcl}
\usepackage[margin=18mm,includeheadfoot,headheight=14pt,headsep=6mm,footskip=9mm]{geometry}
\usepackage{fontspec}
\usepackage{mathtools}
\usepackage{unicode-math}
\setmainfont{TeX Gyre Pagella}
\setsansfont{TeX Gyre Heros}[Scale=MatchLowercase]
\setmonofont{Latin Modern Mono}
\setmathfont{TeX Gyre Pagella Math}
\usepackage[protrusion=true,expansion=false]{microtype}
\usepackage{xcolor}
\usepackage{xurl}
\usepackage[unicode,colorlinks=true,linkcolor=blue!50!black,urlcolor=blue!50!black,
  pdftitle={Sequential Conditional Means of Equicorrelated Scores},
  pdfauthor={Prepared with OpenAI Codex},
  pdfsubject={A fixed starting score included in the average of standardized Gaussian scores}]{hyperref}
\usepackage{bookmark}
\usepackage{scrlayer-scrpage}
\DeclareMathOperator{\E}{E}
\DeclareMathOperator{\Var}{Var}
\DeclareMathOperator{\Cov}{Cov}
\newcommand{\T}{\mathsf{T}}
\newcommand{\given}{\,\vert\,}
\newcommand{\one}{\mathbf{1}}
\setlength{\emergencystretch}{1.5em}
\setkomafont{disposition}{\normalfont\sffamily\bfseries}
\setkomafont{pageheadfoot}{\normalfont\sffamily\small}
\clearpairofpagestyles
\ihead{Sequential conditional means}
\ohead{Derivation and assumptions}
\cfoot*{\pagemark}
\pagestyle{scrheadings}
\title{Sequential Conditional Means of Equicorrelated Scores}
\subtitle{Including the fixed starting score in the average}
\author{}
\date{13 November 2026}

\begin{document}
\maketitle
\thispagestyle{plain.scrheadings}
\vspace{-1em}
\begin{abstract}
Each test has population mean \(0\) and standard deviation \(1\). Fix the
first observed score at \(a\). For jointly Gaussian tests with equal
pairwise correlation \(c\), recursively predict each subsequent test by
conditioning on the starting score and on all earlier predictions.
Every subsequent prediction equals \(ca\). Including \(a\) as the first
value, the average of \(a,e_1,\ldots,e_k\) is
\(a(1+kc)/(k+1)\). We derive the Gaussian conditional mean, solve the recursion,
describe the admissible correlations and singular endpoints, and distinguish
this deterministic construction from actual future scores. A counterexample
shows that means, variances and correlations alone do not establish the result.
\end{abstract}

\noindent\textbf{AI disclosure.} This analysis, its mathematical exposition,
supporting repository documentation and typesetting were prepared with OpenAI
Codex and collaborating AI agents. The requester specified the problem,
clarified the indexing and requested this analysis. Mathematical review and
build checks were performed within the same AI-assisted workflow; no independent
human peer review or new empirical psychometric validation is claimed.

\section{Problem and notation}

Let \(k\geq1\) be an integer and write
\[
  X_0,X_1,\ldots,X_k
\]
for standardized test scores. Throughout this analysis, every test has
population mean zero and standard deviation one:
\begin{equation}
  \E[X_j]=0,\qquad \operatorname{SD}(X_j)=1,
  \qquad \Var(X_j)=1,\qquad 0\leq j\leq k.
  \label{eq:standardization}
\end{equation}
These parameters describe the reference distribution of each test; a fixed
observed score \(a\) need not equal zero. We
relabel the original tests \(A_r(c,0,1)\) so that test \(0\) supplies the
starting observation. For distinct indices,
assume \(\Cov(X_i,X_j)=\operatorname{Corr}(X_i,X_j)=c\).

Use \(a\in\mathbb{R}\) for the \emph{fixed starting score}; it is a scalar,
not a function. For each \(1\leq n\leq k\), let \(g_n\) be the Gaussian conditional
mean function of \(X_n\) given \((X_0,\ldots,X_{n-1})\). Define
\begin{align}
  e_1&=g_1(a)=\E[X_1\given X_0=a],\label{eq:first}\\
  e_n&=g_n(a,e_1,\ldots,e_{n-1})\notag\\
     &=\E[X_n\given X_0=a,X_1=e_1,\ldots,X_{n-1}=e_{n-1}],
       \qquad 2\leq n\leq k.\label{eq:definition}
\end{align}
The quantity of interest is
\begin{equation}
  \boxed{B_k=\frac{a+\sum_{n=1}^{k}e_n}{k+1}.}
  \label{eq:target}
\end{equation}
It contains \(k+1\) values: the starting score \(a\) is the first value,
followed by the \(k\) predictions \(e_1,\ldots,e_k\). Thus \(k\) counts
subsequent predictions, while \(k+1\) counts all values in the average.

For continuous distributions, conditioning on a particular score is not a
ratio of probabilities of point events. A conditional mean is ordinarily
defined only almost everywhere. Here, \(g_n\) means the continuous version
supplied by Gaussian regression, derived below. This convention specifies
the recursion even at individual histories of probability zero. The singular
case \(c=1\) will be treated directly on its support.

\section{The joint model and its correlation domain}

Assume \((X_0,\ldots,X_k)\) is jointly Gaussian, allowing degenerate Gaussian
distributions at the endpoints. Its covariance matrix is
\begin{equation}
  R_{k+1}=(1-c)I_{k+1}+c\one_{k+1}\one_{k+1}^{\T},
  \label{eq:covariance}
\end{equation}
where \(I_d\) is the \(d\)-dimensional identity and \(\one_d\) is a column
of \(d\) ones. Vectors orthogonal to \(\one_{k+1}\) have eigenvalue
\(1-c\); the vector \(\one_{k+1}\) has eigenvalue \(1+kc\). Thus a
covariance matrix for these \(k+1\) scores exists exactly when
\begin{equation}
  \boxed{-\frac1k\leq c\leq1.}\label{eq:domain}
\end{equation}
An arbitrarily long equicorrelated collection requires \(0\leq c\leq1\).
A negative fixed \(c\) cannot satisfy the lower bound for all \(k\).

For \(0\leq c\leq1\), a construction consistent with the hypothetical
common-factor illustration in the repository's
\href{https://github.com/AndersH3/composite_iq_checked_claims/blob/main/analyses/hoeflin-ferguson/composite_iq_analysis.tex}{Hoeflin/Ferguson analysis, Appendix A}, is
\begin{equation}
  X_j=\sqrt c\,G+\sqrt{1-c}\,\varepsilon_j,
  \label{eq:factor}
\end{equation}
where \(G,\varepsilon_0,\ldots,\varepsilon_k\) are independent standard-normal
variables. The covariance-matrix derivation below also covers the admissible
negative correlations. This is a mathematical model, not a claim that actual
test batteries satisfy these assumptions.

\section{Derivation of the conditional mean}

Fix \(1\leq n\leq k\), initially with \(c<1\), and put
\[
  Y_n=(X_0,\ldots,X_{n-1})^{\T},\qquad
  R_n=\Cov(Y_n)=(1-c)I_n+c\one_n\one_n^{\T}.
\]
The eigenvalue of \(R_n\) in the direction \(\one_n\) is
\[
  d_n=1+(n-1)c.
\]
It is positive throughout \(-1/k\leq c<1\): for negative \(c\),
\(d_n\geq1-(n-1)/k\geq1/k>0\). The other eigenvalues are \(1-c>0\).
Consequently every conditioning block is nonsingular, even if the full
\((k+1)\)-score distribution is singular at \(c=-1/k\).

An explicit inverse, verified by multiplying it by \(R_n\), is
\begin{equation}
  R_n^{-1}=\frac1{1-c}
  \left(I_n-\frac{c}{1+(n-1)c}\one_n\one_n^{\T}\right).
  \label{eq:inverse}
\end{equation}
Since \(\Cov(X_n,Y_n)=c\one_n^{\T}\), the regression coefficient is
\begin{equation}
  \Cov(X_n,Y_n)R_n^{-1}
    =\frac{c}{1+(n-1)c}\one_n^{\T}.
  \label{eq:coefficient}
\end{equation}

To establish why this coefficient gives the conditional mean, define the
residual
\begin{equation}
  W_n=X_n-\frac{c}{1+(n-1)c}\sum_{j=0}^{n-1}X_j.
  \label{eq:residual}
\end{equation}
For each \(\ell\in\{0,\ldots,n-1\}\),
\begin{align*}
  \Cov(W_n,X_\ell)
  &=c-\frac{c}{1+(n-1)c}
         \sum_{j=0}^{n-1}\Cov(X_j,X_\ell)\\
  &=c-\frac{c}{1+(n-1)c}\bigl[1+(n-1)c\bigr]=0.
\end{align*}
The residual and \(Y_n\) are jointly Gaussian because they are linear
transformations of a Gaussian vector. Their zero cross-covariance implies
independence: the Gaussian characteristic function
\(\exp(-t^{\T}\Sigma t/2)\) factors across covariance blocks when the
cross-covariance block is zero. This argument also allows a zero-variance
residual. Since \(\E[W_n]=0\), independence gives
\(\E[W_n\given Y_n]=0\). Taking the conditional expectation in
\eqref{eq:residual} therefore yields
\begin{equation}
  \boxed{
    g_n(y_0,\ldots,y_{n-1})
    =\frac{c}{1+(n-1)c}\sum_{j=0}^{n-1}y_j.}
  \label{eq:conditional-mean}
\end{equation}
This linear function is the unique continuous version, because the
nonsingular vector \(Y_n\) has a strictly positive density everywhere.

For completeness, the same residual calculation gives
\begin{align}
  \Var(X_n\given Y_n)
    &=\Var(W_n)
      =1-\frac{nc^2}{1+(n-1)c}\notag\\
    &=\frac{(1-c)(1+nc)}{1+(n-1)c}.
    \label{eq:conditional-variance}
\end{align}
This nonnegative conditional variance describes actual test outcomes,
not randomness in the recursively specified numbers \(e_n\).

\section{Solving the recursion and taking the average}

Inserting the prescribed history into \eqref{eq:conditional-mean} gives
\begin{equation}
  e_n=\frac{c}{1+(n-1)c}
         \left(a+\sum_{j=1}^{n-1}e_j\right),
  \qquad 1\leq n\leq k,
  \label{eq:recursion}
\end{equation}
where the sum is empty at \(n=1\). In particular, \(e_1=ca\).
Suppose \(e_j=ca\) for all \(1\leq j<n\). Then
\[
  e_n=\frac{c}{1+(n-1)c}\bigl[a+(n-1)ca\bigr]
      =\frac{ca[1+(n-1)c]}{1+(n-1)c}=ca.
\]
\begin{samepage}
Induction establishes
\begin{equation}
  \boxed{e_1=e_2=\cdots=e_k=ca.}\label{eq:all-predictions}
\end{equation}
Consequently the requested average is
\begin{equation}
  \boxed{B_k=\frac{a+\sum_{n=1}^{k}e_n}{k+1}
    =\frac{1+kc}{k+1}\,a
    =\left(c+\frac{1-c}{k+1}\right)a.}
  \label{eq:answer}
\end{equation}
\end{samepage}
Unlike the individual subsequent predictions, the requested average depends
on the number of values included. With
\(\lambda_k=(1+kc)/(k+1)\), it is \(B_k=\lambda_k a\).
Once \(a,c,k\) are fixed, this average is deterministic. An additional
expectation changes nothing: \(\E[B_k]=\lambda_k a\).
If \(N=k+1\) denotes the \emph{total} number of values, the same result is
\begin{equation}
  \boxed{B_{N-1}=\frac{1+(N-1)c}{N}\,a,\qquad N\geq2.}
  \label{eq:total-count}
\end{equation}
For the single-value case \(N=1\), the average is simply \(a\), without
any subsequent prediction or inter-test correlation being needed.

There is also a partial-sum derivation. Define
\[
  T_0=a,\qquad T_n=a+\sum_{j=1}^{n}e_j.
\]
Equation~\eqref{eq:recursion} implies
\[
  T_n=\frac{1+nc}{1+(n-1)c}T_{n-1}.
\]
Multiplying these relations telescopes, without dividing by \(c\) or by
\(a\):
\[
  T_n=a\prod_{j=1}^{n}\frac{1+jc}{1+(j-1)c}=a(1+nc).
\]
Dividing \(T_k=a+\sum_{j=1}^{k}e_j=a(1+kc)\) by \(k+1\) again gives
\eqref{eq:answer}. This proof includes \(a=0\) and \(c=0\).

For example, if \(a=2\) and \(c=0.6\), then
\[
  e_1=1.2,\qquad
  e_2=\frac{0.6}{1.6}(2+1.2)=1.2,\qquad
  e_3=\frac{0.6}{2.2}(2+1.2+1.2)=1.2.
\]
Every subsequent prediction remains \(1.2\). Including the starting score,
the average after three predictions is instead
\[
  B_3=\frac{2+1.2+1.2+1.2}{4}
      =\frac{1+3(0.6)}4\,2=1.4.
\]

\begin{samepage}
\section{Interpretation and other averages}

\subsection{Actual scores and their conditional average}
The recursion replaces each previous future score by its conditional mean.
Actual observed future scores need not equal those replacements. If
\[
  C_k=\frac1{k+1}\sum_{j=0}^{k}X_j
\]
is the actual average including the starting test, bivariate Gaussian regression gives
\(\E[X_n\given X_0=a]=ca\) for every \(1\leq n\leq k\). Linearity therefore gives
\begin{equation}
  \E[C_k\given X_0=a]
    =\frac{a+kca}{k+1}=\lambda_k a=B_k.
    \label{eq:actual-mean}
\end{equation}
\end{samepage}
Although this conditional expectation equals \(B_k\), the random
variable \(C_k\) generally has positive conditional variance.

\begin{samepage}
In fact,
\[
  \Var(C_k)=\frac{1+kc}{k+1}=\lambda_k,
  \qquad \Cov(C_k,X_0)=\lambda_k,
\]
so Gaussian regression gives
\begin{equation}
  \Var(C_k\given X_0=a)
  =\lambda_k-\lambda_k^2
  =\frac{k(1-c)(1+kc)}{(k+1)^2}.
  \label{eq:actual-variance}
\end{equation}
\end{samepage}

Likewise, the random predictors
\[
  P_n=\E[X_n\given X_0,\ldots,X_{n-1}]
\]
use actual prior outcomes and are not the deterministic values \(e_n\).
The tower property implies \(\E[P_n\given X_0]=cX_0\), while its
unconditional implication \(\E[P_n]=0\) holds without Gaussianity whenever
the scores are integrable and have mean zero. Applying a conditional mean
function at a fixed, recursively chosen history is a different operation.

\subsection{The subsequent predictions considered separately}
For comparison, averaging just the \(k\) subsequent predictions gives
\begin{equation}
  \frac1k\sum_{n=1}^{k}e_n=ca.
  \label{eq:predictions-only}
\end{equation}
The requested average \(B_k\) includes \(a\) as well and therefore has
denominator \(k+1\), as in \eqref{eq:answer}.

\subsection{Population mean zero and the fixed starting score}
If one randomizes the starting observation according to its standardized
test distribution, the recursively filled average becomes \(\lambda_k X_0\).
Consequently
\begin{equation}
  \E[\lambda_k X_0]=0,\qquad
  \operatorname{SD}(\lambda_k X_0)=\lambda_k,
  \label{eq:random-start}
\end{equation}
where \(\lambda_k\geq0\) throughout the admissible correlation domain.
These are population statements obtained by varying the starting score.
For the fixed starting score in this problem, the answer is
\(B_k=\lambda_k a\), which need not be zero. Every original test still has
mean \(0\) and standard deviation \(1\); the arithmetic average has its own
distribution and is not rescaled to unit standard deviation.

\section{Boundary cases}

At \(c=0\), joint Gaussianity makes the scores independent, and every
subsequent prediction is zero. Including the starting score gives
\(B_k=a/(k+1)\). At \(c=1\),
\(\Var(X_i-X_j)=2(1-c)=0\), so all scores agree almost surely. Given
the starting score \(a\), every compatible recursive prediction is
\(a=ca\), and \(B_k=a\). This establishes the result at \(c=1\) without using the
singular inverse in \eqref{eq:inverse}.

At the other endpoint, \(c=-1/k\), the full vector satisfies
\[
  \sum_{j=0}^{k}X_j=0\quad\text{almost surely},
\]
because that sum has mean zero and variance
\((k+1)(1+kc)=0\). Every conditioning prefix remains nonsingular, as
shown above, so the derivation of \eqref{eq:answer} applies. Its prescribed
history is compatible with the support, since
\(a+kca=0\), so \(B_k=0\). The final conditional variance
\eqref{eq:conditional-variance} is zero.

\section{Why correlation alone is insufficient}

The original description specifies only marginal means, marginal variances
and pairwise correlations. Those quantities do not determine the conditional
mean functions. A counterexample already exists with one subsequent test.
Let
\[
  P(X_0=-\sqrt2)=P(X_0=\sqrt2)=\tfrac14,
  \qquad P(X_0=0)=\tfrac12,
  \qquad X_1=X_0^2-1.
\]
Symmetry and direct calculation give
\[
  \E[X_0]=0,\quad \E[X_0^2]=1,\quad
  \E[X_0^3]=0,\quad \E[X_0^4]=2.
\]
Consequently
\begin{align*}
  \E[X_1]&=\E[X_0^2]-1=0,\\
  \Var(X_1)&=\E[(X_0^2-1)^2]=2-2+1=1,\\
  \Cov(X_0,X_1)&=\E[X_0^3-X_0]=0.
\end{align*}
Both scores are standardized and \(c=0\), yet at the fixed starting value
\(a=0\), an event with positive probability,
\[
  e_1=\E[X_1\given X_0=0]=-1\ne0=ca.
\]
The average including the starting value is therefore
\[
  B_1=\frac{0+(-1)}2=-\frac12,
\]
whereas \eqref{eq:answer} would give zero under the Gaussian model.
Therefore \eqref{eq:answer} is a result under the stated joint Gaussian
model, or another model with the same conditional-mean functions; it does
not follow from equicorrelation alone.

\noindent\textbf{Scope and provenance.} This derivation and its AI review
extend the repository's hypothetical Gaussian model. Existing Ada/SPARK and
Python/Deal evidence concerns the specified software; it does not formally
verify this text. No new empirical study or formal-prover campaign was
performed. The accompanying README records build and visual checks. The
repository's original material follows its existing license, subject to
third-party rights.

\end{document}
